\chapter{Математические модели акустических волноводов}
\label{ch:theory}

В данной главе рассматриваются теоретические аспекты распространения звуковых волн в волноводах переменного сечения в одномерном приближении. Дается вывод уравнения Вебстера и приводится описание метода матриц передачи (TMM) для цилиндрической и конической аппроксимаций геометрии.

\section{Уравнение рупора Вебстера}
\label{sec:webster}

Распространение звуковых волн в жесткостенном канале переменного поперечного сечения $A(x)$ описывается классическим волновым уравнением рупора Вебстера (Webster horn equation)~\cite{webster1919}. Это одномерное уравнение является хорошим приближением для волн, длина которых существенно превышает поперечные размеры канала (то есть для плоских волн).

Математическая формулировка уравнения Вебстера для акустического давления $p(x, t)$ имеет следующий вид:
\begin{equation}
    \frac{1}{A(x)} \frac{\partial}{\partial x} \left( A(x) \frac{\partial p(x, t)}{\partial x} \right) - \frac{1}{c^2} \frac{\partial^2 p(x, t)}{\partial t^2} = 0,
    \label{eq:webster_time}
\end{equation}
где $x$ --- координата вдоль оси волновода, $t$ --- время, $c$ --- скорость звука в среде.

Для гармонических волн с временной зависимостью $e^{j\omega t}$ уравнение Вебстера во временной области переходит в уравнение в частотной области:
\begin{equation}
    \frac{d^2 P(x, \omega)}{dx^2} + \frac{1}{A(x)} \frac{dA(x)}{dx} \frac{dP(x, \omega)}{dx} + k^2 P(x, \omega) = 0,
    \label{eq:webster_freq}
\end{equation}
где $P(x, \omega)$ --- комплексная амплитуда акустического давления, $k = \omega / c$ --- волновое число, $\omega = 2\pi f$ --- круговая частота.

Если ввести вспомогательную функцию $\psi(x) = P(x) \sqrt{A(x)}$, то уравнение~\eqref{eq:webster_freq} может быть преобразовано к одномерному уравнению Шредингера:
\begin{equation}
    \frac{d^2 \psi(x)}{dx^2} + [k^2 - V(x)] \psi(x) = 0,
\end{equation}
где $V(x) = \frac{1}{2 A} \frac{d^2 A}{dx^2} - \frac{1}{4 A^2} \left(\frac{dA}{dx}\right)^2$ представляет собой так называемый «потенциал рупора».

\section{Метод матриц передачи (TMM)}
\label{sec:tmm}

Метод матриц передачи (Transfer Matrix Method, TMM) основан на представлении сложного акустического волновода в виде каскадного соединения более простых элементов (сегментов), для каждого из которых известно точное аналитическое решение. Переменными состояния на границах сегментов являются акустическое давление $P$ и объемная колебательная скорость $U$.

Для любого акустического элемента связь между входными ($P_{in}, U_{in}$) и выходными ($P_{out}, U_{out}$) величинами записывается через матрицу передачи $T$ размера $2 \times 2$:
\begin{equation}
    \begin{pmatrix} P_{in} \\ U_{in} \end{pmatrix} = 
    \begin{pmatrix} T_{11} & T_{12} \\ T_{21} & T_{22} \end{pmatrix}
    \begin{pmatrix} P_{out} \\ U_{out} \end{pmatrix}.
\end{equation}

При каскадном соединении $N$ сегментов общая матрица передачи всей системы равна произведению матриц передачи отдельных сегментов:
\begin{equation}
    T_{total} = T_1 \cdot T_2 \cdot \dots \cdot T_N.
\end{equation}

\subsection{Цилиндрическая аппроксимация геометрии}
Наиболее простым способом дискретизации является разбиение волновода на $N$ цилиндрических сегментов постоянного сечения. Для цилиндра длиной $l$ и площадью поперечного сечения $A$ элементы матрицы передачи имеют вид~\cite{chiba2006}:
\begin{equation}
    T_{cyl} = \begin{pmatrix} 
        \cos(kl) & j Z_0 \sin(kl) \\ 
        \frac{j}{Z_0} \sin(kl) & \cos(kl) 
    \end{pmatrix},
\end{equation}
где $Z_0 = \frac{\rho_0 c}{A}$ --- характеристический волновой импеданс цилиндрического волновода, $\rho_0$ --- плотность среды.

\subsection{Коническая аппроксимация геометрии}
Коническая аппроксимация волновода обеспечивает более точное представление геометрии с непрерывным изменением площади. Для конического сегмента с площадями на входе $A_{in}$ (радиус $r_{in}$) и выходе $A_{out}$ (радиус $r_{out}$) и длиной $l$ элементы матрицы передачи выражаются через сферические функции:
\begin{equation}
    T_{cone} = \begin{pmatrix}
        \frac{r_{out}}{r_{in}} \cos(kl) - \frac{1}{k r_{in}} \sin(kl) & j \frac{\rho_0 c}{S_{in}} \frac{r_{out}}{r_{in}} \sin(kl) \\
        \dots & \dots
    \end{pmatrix}.
\end{equation}
Использование конических сегментов позволяет снизить число необходимых элементов разбиения для достижения той же точности по сравнению с чисто цилиндрической аппроксимацией.
